% GENERATED FILE. Edit canonical lessons, then run npm run generate:books.
% canonical-source-hash: fnv1a64:16f0e2a3707bfbf0
% canonical-lessons: BN-C399-bhromon-kora, BN-C399-bonbhojon, BN-C399-tabu, BN-C399-pohano, BN-C399-gaid, BN-R399-review-392-399

\chapter{Trips and Camps}
\label{ch:bn-trips-and-camps}
\hlchaptermodality{399}

\begin{chapteropening}
\textbf{By the end of this chapter:} \emph{I can talk about travelling, a picnic and a night in a tent by the fire.}

Complete the last lesson of chapter 399: I can talk about travelling, a picnic and a night in a tent by the fire.
\end{chapteropening}

\section[\texorpdfstring{\bn{ভ্রমণ} \bn{করা} (bhrômôn kôrā)}{ভ্রমণ করা (bhrômôn kôrā)}]{\bn{ভ্রমণ} \bn{করা} (bhrômôn kôrā) — to travel; to tour}

\label{lesson:BN-C399-bhromon-kora}



\begin{quote}
Before the new one: say the Bengali for to take a dip, then the Bengali for to splash.
\end{quote}

\subsection*{You'll want to know: \bn{ভ্রমণ} \bn{করা}}
\textbf{\bn{ভ্রমণ} \bn{করা}} — \emph{bhrômôn kôrā} — \textquotedblleft{}to travel; to tour\textquotedblright{}.

\begin{quote}
\bn{গত} \bn{বছর} \bn{আমরা} \bn{পাহাড়ে} \bn{ভ্রমণ} \bn{করলাম।}
\end{quote}

\begin{quote}
\emph{gôtô bôchhôr āmrā pāhāṛe bhrômôn kôrlām.}
\end{quote}

\begin{quote}
Last year we travelled in the hills.
\end{quote}

Past (he/she …): \bn{তিনি} \bn{অনেক} \bn{দেশ} \bn{ভ্রমণ} \bn{করলেন}, \emph{tini ônek desh bhrômôn kôrlen}.

A careful word, for travel writing and the news; \bn{বেড়ানো} is the everyday one.

\begin{cousinweb}
\textbf{\bn{ভ্রমণ} \bn{করা}} (\emph{bhrômôn kôrā}) joins Sanskrit \emph{bhramaṇa}, \textquotedblleft{}wandering, roaming about\textquotedblright{} (Monier-Williams), to \bn{করা}, to do.
\end{cousinweb}

\begin{grammarlens}[title={What you've built}]
One more word for getting around.
\end{grammarlens}

\subsection*{Your turn}
\begin{itemize}
\raggedright
  \item \emph{Say it:} \emph{bhrômôn kôrā}
  \item \emph{Say it:} \emph{bhrômôn kôrā}, once more
  \item \emph{Say it:} \emph{gôtô bôchhôr āmrā pāhāṛe bhrômôn kôrlām.} — Last year we travelled in the hills.
  \item \emph{Recall:} say the Bengali for to take a dip, then the Bengali for to splash, then say \emph{bhrômôn kôrā} again
\end{itemize}

\begin{tcolorbox}[breakable,colback=teal!4,colframe=teal!35!black,title={Before you move on}]
What does \bn{ভ্রমণ} \bn{করা} mean? (\textquotedblleft{}To travel; to tour\textquotedblright{}.) Say it once more.
\end{tcolorbox}
\section[\texorpdfstring{\bn{বনভোজন} (bônbhojôn)}{বনভোজন (bônbhojôn)}]{\bn{বনভোজন} (bônbhojôn) — a picnic}

\label{lesson:BN-C399-bonbhojon}



\begin{quote}
Before the new one: say the Bengali for to splash, then the Bengali for to travel.
\end{quote}

\subsection*{You'll want to know: \bn{বনভোজন}}
\textbf{\bn{বনভোজন}} — \emph{bônbhojôn} — \textquotedblleft{}a picnic\textquotedblright{}.

\begin{quote}
\bn{শীতে} \bn{আমরা} \bn{নদীর} \bn{পাড়ে} \bn{বনভোজন} \bn{করলাম।}
\end{quote}

\begin{quote}
\emph{shīte āmrā nôdīr pāṛe bônbhojôn kôrlām.}
\end{quote}

\begin{quote}
In winter we had a picnic on the riverbank.
\end{quote}

Winter is the season for it in Bengal, with a cook-pot carried out to a field or a riverbank.

\begin{cousinweb}
\textbf{\bn{বনভোজন}} (\emph{bônbhojôn}) joins \textbf{\bn{বন}}, \textquotedblleft{}a forest\textquotedblright{}, Sanskrit \emph{vana}, to Sanskrit \emph{bhojana}, \textquotedblleft{}eating, a meal\textquotedblright{} (Monier-Williams for both): a meal out in the woods.
\end{cousinweb}

\begin{grammarlens}[title={What you've built}]
One more word for getting around.
\end{grammarlens}

\subsection*{Your turn}
\begin{itemize}
\raggedright
  \item \emph{Say it:} \emph{bônbhojôn}
  \item \emph{Say it:} \emph{bônbhojôn}, once more
  \item \emph{Say it:} \emph{shīte āmrā nôdīr pāṛe bônbhojôn kôrlām.} — In winter we had a picnic on the riverbank.
  \item \emph{Recall:} say the Bengali for to splash, then the Bengali for to travel, then say \emph{bônbhojôn} again
\end{itemize}

\begin{tcolorbox}[breakable,colback=teal!4,colframe=teal!35!black,title={Before you move on}]
What does \bn{বনভোজন} mean? (\textquotedblleft{}A picnic\textquotedblright{}.) Say it once more.
\end{tcolorbox}
\section[\texorpdfstring{\bn{তাঁবু} (tā̃bu)}{তাঁবু (tā̃bu)}]{\bn{তাঁবু} (tā̃bu) — a tent}

\label{lesson:BN-C399-tabu}



\begin{quote}
Before the new one: say the Bengali for to travel, then the Bengali for a picnic.
\end{quote}

\subsection*{You'll want to know: \bn{তাঁবু}}
\textbf{\bn{তাঁবু}} — \emph{tā̃bu} — \textquotedblleft{}a tent\textquotedblright{}.

\begin{quote}
\bn{পাহাড়ে} \bn{আমরা} \bn{তাঁবুতে} \bn{রাত} \bn{কাটালাম।}
\end{quote}

\begin{quote}
\emph{pāhāṛe āmrā tā̃bute rāt kāṭālām.}
\end{quote}

\begin{quote}
In the hills we spent the night in a tent.
\end{quote}

\begin{cousinweb}
\textbf{\bn{তাঁবু}} (\emph{tā̃bu}) is usually traced to Persian \emph{tanbū}, \textquotedblleft{}a tent\textquotedblright{}, with the n become the nasal \bn{ঁ}.
\end{cousinweb}

\begin{grammarlens}[title={What you've built}]
One more word for getting around.
\end{grammarlens}

\subsection*{Your turn}
\begin{itemize}
\raggedright
  \item \emph{Say it:} \emph{tā̃bu}
  \item \emph{Say it:} \emph{tā̃bu}, once more
  \item \emph{Say it:} \emph{pāhāṛe āmrā tā̃bute rāt kāṭālām.} — In the hills we spent the night in a tent.
  \item \emph{Recall:} say the Bengali for to travel, then the Bengali for a picnic, then say \emph{tā̃bu} again
\end{itemize}

\begin{tcolorbox}[breakable,colback=teal!4,colframe=teal!35!black,title={Before you move on}]
What does \bn{তাঁবু} mean? (\textquotedblleft{}A tent\textquotedblright{}.) Say it once more.
\end{tcolorbox}
\section[\texorpdfstring{\bn{পোহানো} (pohāno)}{পোহানো (pohāno)}]{\bn{পোহানো} (pohāno) — to warm oneself (at a fire or in the sun)}

\label{lesson:BN-C399-pohano}



\begin{quote}
Before the new one: say the Bengali for a picnic, then the Bengali for a tent.
\end{quote}

\subsection*{You'll want to know: \bn{পোহানো}}
\textbf{\bn{পোহানো}} — \emph{pohāno} — \textquotedblleft{}to warm oneself (at a fire or in the sun)\textquotedblright{}.

\begin{quote}
\bn{রাতে} \bn{তাঁবুর} \bn{বাইরে} \bn{আমরা} \bn{আগুন} \bn{পোহালাম।}
\end{quote}

\begin{quote}
\emph{rāte tā̃bur bāire āmrā āgun pohālām.}
\end{quote}

\begin{quote}
At night we warmed ourselves at a fire outside the tent.
\end{quote}

Past (he/she …): \bn{তিনি} \bn{আগুন} \bn{পোহালেন}, \emph{tini āgun pohālen}.

\bn{আগুন} \bn{পোহানো}, to warm oneself at a fire; \bn{রোদ} \bn{পোহানো}, to sit in the winter sun.

\begin{cousinweb}
\textbf{\bn{পোহানো}} (\emph{pohāno}) is inherited: CDIAL lists Bengali \emph{pohāna}, \textquotedblleft{}to warm oneself by the fire\textquotedblright{}, under Sanskrit \emph{pratāpayati}, \textquotedblleft{}warms, heats\textquotedblright{} (CDIAL 8539).
\end{cousinweb}

\begin{grammarlens}[title={What you've built}]
One more word for getting around.
\end{grammarlens}

\subsection*{Your turn}
\begin{itemize}
\raggedright
  \item \emph{Say it:} \emph{pohāno}
  \item \emph{Say it:} \emph{pohāno}, once more
  \item \emph{Say it:} \emph{rāte tā̃bur bāire āmrā āgun pohālām.} — At night we warmed ourselves at a fire outside the tent.
  \item \emph{Recall:} say the Bengali for a picnic, then the Bengali for a tent, then say \emph{pohāno} again
\end{itemize}

\begin{tcolorbox}[breakable,colback=teal!4,colframe=teal!35!black,title={Before you move on}]
What does \bn{পোহানো} mean? (\textquotedblleft{}To warm oneself (at a fire or in the sun)\textquotedblright{}.) Say it once more.
\end{tcolorbox}
\section[\texorpdfstring{\bn{গাইড} (gāiḍ)}{গাইড (gāiḍ)}]{\bn{গাইড} (gāiḍ) — a guide (for visitors)}

\label{lesson:BN-C399-gaid}



\begin{quote}
Before the new one: say the Bengali for a tent, then the Bengali for to warm oneself.
\end{quote}

\subsection*{You'll want to know: \bn{গাইড}}
\textbf{\bn{গাইড}} — \emph{gāiḍ} — \textquotedblleft{}a guide (for visitors)\textquotedblright{}.

\begin{quote}
\bn{গাইড} \bn{আমাদের} \bn{পুরনো} \bn{দুর্গটা} \bn{দেখালেন।}
\end{quote}

\begin{quote}
\emph{gāiḍ āmāder purôno durgôṭā dekhālen.}
\end{quote}

\begin{quote}
The guide showed us the old fort.
\end{quote}

\begin{cousinweb}
\textbf{\bn{গাইড}} (\emph{gāiḍ}) is borrowed from English \emph{guide}, with the final d said as the hard \bn{ড}.
\end{cousinweb}

\begin{grammarlens}[title={What you've built}]
One more word for getting around.
\end{grammarlens}

\subsection*{Your turn}
\begin{itemize}
\raggedright
  \item \emph{Say it:} \emph{gāiḍ}
  \item \emph{Say it:} \emph{gāiḍ}, once more
  \item \emph{Say it:} \emph{gāiḍ āmāder purôno durgôṭā dekhālen.} — The guide showed us the old fort.
  \item \emph{Recall:} say the Bengali for a tent, then the Bengali for to warm oneself, then say \emph{gāiḍ} again
\end{itemize}

\begin{tcolorbox}[breakable,colback=teal!4,colframe=teal!35!black,title={Before you move on}]
What does \bn{গাইড} mean? (\textquotedblleft{}A guide (for visitors)\textquotedblright{}.) Say it once more.
\end{tcolorbox}
\section[(dialogue)]{Looking back: from Traffic and Crossings to Trips and Camps}

\label{lesson:BN-R399-review-392-399}



\begin{quote}
Say the Bengali for traffic and a pedestrian.
\end{quote}

\subsection*{Your turn}
Say these aloud:

\begin{itemize}
\raggedright
  \item traffic, a pedestrian, a flyover, a signboard, to keep close to
  \item diesel, air, a puncture, to run out, to start (an engine)
  \item to break down, a spare part, a screw, a steering wheel, to rub
  \item a flight, a runway, customs, abroad, to make a crossing
  \item a boatman, an oar, an anchor, the hood of a country boat, to row
  \item a landing place, a launch, a ferry, a bank (of a pond or river), to moor
  \item a current, a beach, a jetty, to take a dip, to splash
  \item to travel, a picnic, a tent, to warm oneself, a guide (for visitors)
\end{itemize}

\begin{tcolorbox}[breakable,colback=teal!4,colframe=teal!35!black,title={Before you move on}]
Traffic? (\textbf{\bn{ট্রাফিক}}, \emph{ṭrāphik}.) A pedestrian? (\textbf{\bn{পথচারী}}, \emph{pôthchārī}.) A flyover? (\textbf{\bn{ফ্লাইওভার}}, \emph{phlāiobhār}.) A signboard? (\textbf{\bn{সাইনবোর্ড}}, \emph{sāinborḍ}.) To keep close to? (\textbf{\bn{ঘেঁষা}}, \emph{ghẽshā}.) Diesel? (\textbf{\bn{ডিজেল}}, \emph{ḍijel}.) Air? (\textbf{\bn{হাওয়া}}, \emph{hāoyā}.) A puncture? (\textbf{\bn{পাংচার}}, \emph{pāngchār}.) To run out? (\textbf{\bn{ফুরানো}}, \emph{phurāno}.) To start (an engine)? (\textbf{\bn{চালু} \bn{করা}}, \emph{chālu kôrā}.) To break down? (\textbf{\bn{বিগড়ানো}}, \emph{bigṛāno}.) A spare part? (\textbf{\bn{যন্ত্রাংশ}}, \emph{jôntrāngshô}.) A screw? (\textbf{\bn{স্ক্রু}}, \emph{skru}.) A steering wheel? (\textbf{\bn{স্টিয়ারিং}}, \emph{sṭiyāring}.) To rub? (\textbf{\bn{ঘষা}}, \emph{ghôshā}.) A flight? (\textbf{\bn{ফ্লাইট}}, \emph{phlāiṭ}.) A runway? (\textbf{\bn{রানওয়ে}}, \emph{rānoye}.) Customs? (\textbf{\bn{কাস্টমস}}, \emph{kāsṭôms}.) Abroad? (\textbf{\bn{বিদেশ}}, \emph{bidesh}.) To make a crossing? (\textbf{\bn{পাড়ি} \bn{দেওয়া}}, \emph{pāṛi deoyā}.) A boatman? (\textbf{\bn{মাঝি}}, \emph{mājhi}.) An oar? (\textbf{\bn{দাঁড়}}, \emph{dā̃ṛ}.) An anchor? (\textbf{\bn{নোঙর}}, \emph{nongôr}.) The hood of a country boat? (\textbf{\bn{ছই}}, \emph{chhôi}.) To row? (\textbf{\bn{বাওয়া}}, \emph{bāoyā}.) A landing place? (\textbf{\bn{ঘাট}}, \emph{ghāṭ}.) A launch? (\textbf{\bn{লঞ্চ}}, \emph{lônch}.) A ferry? (\textbf{\bn{খেয়া}}, \emph{kheyā}.) A bank (of a pond or river)? (\textbf{\bn{পাড়}}, \emph{pāṛ}.) To moor? (\textbf{\bn{ভেড়ানো}}, \emph{bheṛāno}.) A current? (\textbf{\bn{স্রোত}}, \emph{srot}.) A beach? (\textbf{\bn{সৈকত}}, \emph{soikôt}.) A jetty? (\textbf{\bn{জেটি}}, \emph{jeṭi}.) To take a dip? (\textbf{\bn{ডুব} \bn{দেওয়া}}, \emph{ḍub deoyā}.) To splash? (\textbf{\bn{ছিটানো}}, \emph{chhiṭāno}.) To travel? (\textbf{\bn{ভ্রমণ} \bn{করা}}, \emph{bhrômôn kôrā}.) A picnic? (\textbf{\bn{বনভোজন}}, \emph{bônbhojôn}.) A tent? (\textbf{\bn{তাঁবু}}, \emph{tā̃bu}.) To warm oneself? (\textbf{\bn{পোহানো}}, \emph{pohāno}.) A guide (for visitors)? (\textbf{\bn{গাইড}}, \emph{gāiḍ}.)
\end{tcolorbox}
